\documentclass[conference]{IEEEtran}
\IEEEoverridecommandlockouts
% The preceding line is only needed to identify funding in the first footnote. If that is unneeded, please comment it out.
\usepackage{cite}
\usepackage{amsmath,amssymb,amsfonts}
\usepackage{algorithmic}
\usepackage{graphicx}
\graphicspath{{./}}
\usepackage{textcomp}
\usepackage{xcolor}
\usepackage{url}
\usepackage{booktabs}
\usepackage{multirow}
\usepackage{array}
\usepackage{colortbl}
\usepackage{makecell}
\def\BibTeX{{\rm B\kern-.05em{\sc i\kern-.025em b}\kern-.08em
    T\kern-.1667em\lower.7ex\hbox{E}\kern-.125emX}}
\begin{document}

\title{GNN\_KAN: Diagnosing Non-Linear Fault Propagation in Scale-Free Microservices via Kolmogorov-Arnold Networks}

% \author{\IEEEauthorblockN{1\textsuperscript{st} Kuan-Hao Huang}
% \IEEEauthorblockA{\textit{Department of Computer Science} \\
% \textit{National Taichung University of Education}\\
% Taichung, Taiwan \\
% bcs113116@gm.ntcu.edu.tw}
% \and
% \IEEEauthorblockN{2\textsuperscript{nd} Yu-Chen Lai}
% \IEEEauthorblockA{\textit{Department of Computer Science} \\
% \textit{National Taichung University of Education}\\
% Taichung, Taiwan \\
% bcs113115@gm.ntcu.edu.tw}
% \and
% \IEEEauthorblockN{3\textsuperscript{rd} Kuan-Chou Lai}
% \IEEEauthorblockA{\textit{Department of Computer Science} \\
% \textit{National Taichung University of Education}\\
% Taichung, Taiwan \\
% kclai@mail.ntcu.edu.tw}
% }

\maketitle

\begin{abstract}
Root cause analysis (RCA) is the process of identifying the fundamental source of system failures or anomalies. In cloud-native microservices, where applications are decomposed into numerous independent services communicating over networks, rapid fault localization is critical for maintaining system reliability. The mean time to resolution (MTTR) directly impacts operational costs, and delayed diagnosis can lead to cascading failures across service dependencies. However, existing approaches struggle with the complexity of scale-free service topologies and non-linear fault propagation patterns. Standard Graph Attention Networks (GATs) suffer from structural attention bias in scale-free microservice topologies, where scalar attention weights ($\alpha_{ij} \in \mathbb{R}$) correlate with node degree rather than causal dependency, causing attention collapse to hub nodes. This study introduces \textbf{GNN\_KAN} to decouple structural bias from causal learning by replacing scalar weights with learnable vector-valued Kolmogorov-Arnold representations. Furthermore, to capture non-linear fault dynamics, this study incorporates Rotation-based Periodic Kernels (RPK) that model resource saturation thresholds through bounded tensor-product mappings. On the Train-Ticket benchmark (50+ services), GNN\_KAN achieves Precision@1 of 0.272, outperforming GAT (0.064), BARO (0.136), and GNN with MLP (0.024), representing a 1033\% improvement over MLP baselines. The proposed method maintains linear computational complexity $O(|E|)$ with inference latency of 433.2\,ms per snapshot (batch size=1, measured on NVIDIA 2080ti GPU), reducing the diagnostic search space by 95\% from 50+ candidates to Top-5. This study aims to improve root cause localization accuracy while providing interpretable evidence through learnable function shapes that reflect physical resource constraints.
\end{abstract}

\begin{IEEEkeywords}
Root cause analysis, Graph neural networks, Kolmogorov-Arnold Networks, Microservices systems, Resource-constrained message passing
\end{IEEEkeywords}

\section{Introduction}

\subsection{Research Background and Motivation}

Microservices architecture decomposes monolithic applications into small, independently deployable services that communicate via well-defined APIs. This architectural pattern enables scalability and fault isolation, but introduces operational complexity: service dependencies form dynamic graphs where faults propagate across service boundaries. As distributed microservices become the de facto standard for cloud-native applications \cite{microservices_survey}, their scale-free topology creates a degree centrality bias in graph learning models. Scalar attention mechanisms suffer from scalar-to-degree correlation, where attention weights ($\alpha_{ij} \in \mathbb{R}$) statistically align with node degree rather than fault propagation paths. Furthermore, system faults often stem from resource saturation, which exhibits non-linear asymptotic behaviors (e.g., performance collapse after 85\% memory usage). Standard GNNs with unbounded activations (e.g., ReLU) struggle to model soft-saturation limits. To address this limitation, this study introduces a resource-constrained graph neural network message passing mechanism, leveraging the Kolmogorov--Arnold representation theorem \cite{kolmogorov1957,kan2024} to replace scalar weights with learnable univariate basis functions that model differentiable saturation boundaries on edges.

\textit{Diagnostic Efficiency.} The proposed method achieves a Top-5 Recall of 45.6\% on the Train-Ticket benchmark, effectively isolating root causes within the top 5 candidates in topologies of approximately 50 nodes.

\subsection{Research Contributions}

\begin{itemize}
    \item \textit{KAN-based Graph Neural Network Message Passing:} This study proposes a learnable edge-weighting architecture that substitutes scalar attention with vector-valued KAN layers in the graph neural network message passing framework. This explicitly models the non-linear gating effect of resource constraints on fault propagation, using separable basis functions to approximate physical resource limits.
    
    \item \textit{Rotation-based Periodic Kernel (RPK):} To capture cyclic anomalies (e.g., retry storms) that polynomial bases fail to represent, this study introduces RPK. Unlike standard B-splines used in KANs, RPK utilizes strictly bounded periodic functions within $[-1, 1]$. This formulation enables the modeling of repetitive resource contention patterns while ensuring numerical stability, achieving linear complexity $O(|E|)$ through efficient dense matrix operations (16D feature interaction mapping) with an inference latency of 433.2\,ms per snapshot.
    
    \item \textit{Dual-Frequency Filtering Strategy:} This study designs a heterogeneous basis ensemble to decouple fault dynamics. RPK leverages its strict boundedness to model low-frequency steady-state fluctuations and saturation limits, while Chebyshev polynomials function as a high-frequency filter for capturing transient spikes and rapid trend shifts. An element-wise maximization operator acts as a non-parametric selector, dynamically switching between these frequency regimes based on the dominant fault characteristic.
\end{itemize}

\section{Related Work}

Root cause analysis is a key problem in distributed system reliability engineering. This study categorizes existing approaches into three paradigms following established taxonomies in the literature \cite{microrank2021,eadro2023}: statistical and random walk-based methods, deep learning and causal inference approaches, and knowledge-guided learning.

\subsection{Graph-based Statistical and Random Walk Methods}

Graph-based statistical methods, such as MicroRank \cite{microrank2021}, T-Rank \cite{trank2021}, CloudRanger \cite{cloudranger2018}, and BARO \cite{baro2023}, dominate the landscape of microservice root cause analysis. These approaches typically rely on random walks, spectral analysis, or correlation metrics to rank candidate nodes. However, they rely on the linear propagation assumption: random walk algorithms assume conservation of flow (probability sums to 1), treating fault propagation as a conserved quantity. This constraint prevents them from capturing resource saturation effects, where faults behave as sinks (e.g., memory leaks) or exhibit threshold-based blocking, rather than linear transmission. Even multi-modal approaches like Nezha \cite{nezha2023} ultimately depend on random walk-based reasoning engines, inheriting these linearity constraints.

\subsection{Deep Learning and Causal Inference}

To address non-linearity, deep learning approaches (e.g., EADRO \cite{eadro2023}) employ neural networks for fault modeling. However, standard GNNs predominantly utilize Multi-Layer Perceptrons (MLPs) with piecewise linear activations (ReLU/GeLU). As discussed in Section I, these unbounded activation functions struggle to model the asymptotic boundaries of resource exhaustion.

In the meantime, Causal Inference methods (e.g., \cite{gnn_rca2023}, CausalRCA \cite{causalrca2023}) attempt to reconstruct the underlying causal graph directly. While theoretically sound, they face three practical hurdles in production environments: (1) prohibitive computational costs for structure learning ($O(N!)$ or polynomial with high constants), (2) the requirement for large-sample interventions, and (3) the strict Directed Acyclic Graph (DAG) assumption, which contradicts the cyclic nature of service dependencies (e.g., mutual RPC calls) in real-world microservices.

\subsection{LLMs and KAN-based Approaches}

Recent studies, such as RCACopilot \cite{rcacopilot2024} and Xpert \cite{xpert2024}, leverage Large Language Models (LLMs) for incident diagnosis. While LLMs demonstrate superior capabilities in semantic log parsing, their application to high-frequency numerical metrics incurs prohibitive token costs and inference latency ($O(N^2)$ attention complexity), making them unsuitable for real-time snapshot analysis.

This study integrates Kolmogorov-Arnold Networks (KAN) \cite{kan2024} into the graph message-passing mechanism. KAN, based on the Kolmogorov--Arnold representation theorem \cite{kolmogorov1957,kan2024}, decomposes multivariate functions into learnable, interpretable one-dimensional basis functions. Unlike the black-box transformations of MLPs, KAN enables explicit embedding of resource constraints and saturation limits through bounded basis function shapes.

\section{Methodology}

This study formalizes Root Cause Analysis (RCA) as a learning-to-rank problem on dynamic graphs, where the primary objective is to identify the source node $v^*$ that maximizes causal probability under structural constraints.

\subsection{Problem Definition and Challenges}

Root Cause Analysis (RCA) aims to identify the most likely root cause that triggers an anomaly from a candidate set of service nodes when a distributed system experiences an anomalous event. Given a microservices system graph $G=(V, E)$, where $V$ is the set of service nodes and $E \subseteq V \times V$ is the set of dependency edges. Let $X = \{x_v(t) | v \in V, t \in W\}$ be node observation metrics (such as latency, error rate, throughput, CPU/memory usage) within time window $W$, and $A$ be the anomaly event identifier. The research goal is to learn a ranking function $f: (G, X) \to R$ that computes a root cause score $s_v$ for each node $v$, aiming to rank the true root cause $v^*$ within the top-$k$ candidates in the ranking list $R$.

Core technical challenges include: (1) \textit{Cascade masking effect}: anomalies from root cause nodes propagate downstream, causing downstream symptoms to overshadow upstream root causes; (2) \textit{Dynamic topology}: service dependency graphs change with version updates and dynamic scaling; (3) \textit{Interpretability}: Operational transparency requires the model to provide interpretable evidence for fault determination. GNN\_KAN addresses this by controlling the propagation transformation on every edge with a single KAN function, where the learned function shape reflects specific fault patterns (e.g., saturation).

\subsection{Leakage-Free Feature Construction}

A common challenge in RCA is \textit{posterior leakage}---using statistics (e.g., global min/max) derived from the entire event window to normalize data, inadvertently leaking future information. Standard Z-score normalization uses the global mean of the \textit{entire} test set. This study defines the feature window $W$ as $[t-10, t]$ and normalizes using historical statistics derived from $W_{\text{hist}}$. This study extracts a concise 48-dimensional vector $\bar{x}_v$ focusing on gradients and distributions rather than raw values. Key features include: (1) \textit{Lagged Correlation}: Capturing propagation delays between service-level indicators (latency P50, P95, P99 percentiles; error rate; throughput); (2) \textit{Distributional Shifts}: Skewness and kurtosis changes before/after the trigger, which distinguish sudden crashes (network cuts) from progressive degradation (memory leaks); (3) \textit{Anomaly Intensity}: Peak anomaly intensity and sustained anomaly duration, computed using z-score thresholds; (4) \textit{Temporal Gradients}: Statistical changes (mean, variance) of latency and resource metrics, capturing different fault characteristics. All features are computed within the event window $W$ and historical window $W_{\text{hist}}$ (5 minutes each), excluding future information. Features are standardized using z-score normalization and NaN/Inf values are handled to ensure numerical stability. This framework prevents the model from learning training-set-specific spurious correlations.

\subsection{Graph Construction}

This study represents microservices systems within anomaly event windows using two complementary graph structures: the propagation graph $G_{\text{prop}} = (V, E_{\text{prop}}, W_{\text{prop}})$ and the similarity graph $G_{\text{sim}} = (V, E_{\text{sim}}, W_{\text{sim}})$.

\textit{Propagation Graph ($G_{\text{prop}}$):} The propagation graph $G_{\text{prop}}$ captures structural dependencies derived from service call traces and dependency relationships. The edge set $E_{\text{prop}} \subseteq V \times V$ is extracted from service dependency relationships and call tracing information. The edge weight matrix $W_{\text{prop}} \in \mathbb{R}^{|V| \times |V|}$ is initially normalized based on inter-service call frequencies (dividing call counts by maximum call counts to ensure weights fall in $[0, 1]$), and adaptively adjusted through learnable parameters during training to reflect actual anomaly propagation strength. This initialization strategy uses reasonable structural priors, accelerating convergence.

\textit{Similarity Graph ($G_{\text{sim}}$):} The similarity graph $G_{\text{sim}}$ captures feature-level correlations between nodes based on their anomaly observation patterns. For node feature vectors $\bar{x}_{v_i}, \bar{x}_{v_j} \in \mathbb{R}^{d_{\text{feat}}}$, this study computes cosine similarity \cite{microrank2021}:
\begin{equation}
\text{sim}(v_i, v_j) = \frac{\bar{x}_{v_i} \cdot \bar{x}_{v_j}}{||\bar{x}_{v_i}|| \cdot ||\bar{x}_{v_j}||}
\end{equation}
An edge $(v_i, v_j) \in E_{\text{sim}}$ exists if $\text{sim}(v_i, v_j) > \tau_{\text{sim}}$, where $\tau_{\text{sim}}$ is a hyperparameter determined by sensitivity analysis (set to 0.5 in the experiments). The edge weight $W_{\text{sim}}[i,j] = \text{sim}(v_i, v_j)$ reflects the strength of feature correlation. The similarity graph identifies nodes with similar anomaly signatures, complementing the structural dependencies captured by the propagation graph.

For node $v_i \in V$, this study converts its time-series data within the observation window into a fixed-dimensional feature vector $\bar{x}_{v_i} \in \mathbb{R}^{d_{\text{feat}}}$ through feature extraction, then obtains initial representation $h_{v_i}^{(0)} \in \mathbb{R}^{d_h}$ through linear transformation and LayerNorm, where $d_h$ is the hidden representation dimension (64).

\begin{figure*}[t!]
    \centering
    \includegraphics[width=0.95\textwidth]{figure1_structure.jpg}
    \caption{GNN\_KAN Architecture Overview. The left side shows the dual-graph structure (similarity graph $G_{\text{sim}}$ and propagation graph $G_{\text{prop}}$) constructed based on RCA-aware features; the middle layer shows the message passing mechanism based on Kolmogorov--Arnold representation, replacing traditional MLPs with learnable one-dimensional functions (wavy line symbols), including three core modules: $\text{KAN}_{\text{edge}}$, $\text{KAN}_{\text{message}}$, and $\text{KAN}_{\text{node}}$; the right side shows the $\text{KAN}_{\text{fusion}}$ fusion layer and readout layer, replacing opaque MLP transformations with interpretable functional mappings. The figure labels two main paths: the similarity path (processing similarity features $h^{\text{sim}}$) and the propagation path (processing propagation features $h^{\text{prop}}$), ultimately generating root cause scores $s_v$ through the fusion layer.}
    \label{fig:architecture}
\end{figure*}

\subsection{GNN\_KAN Model Overview}

GNN\_KAN integrates Kolmogorov--Arnold representation networks into graph neural network message passing. The overall model and processing flow are shown in Fig.~\ref{fig:architecture}.

This study represents microservices systems as dynamic weighted directed graphs $G = (V, E, w)$, where $w: E \to \mathbb{R}$ is an edge weight function reflecting anomaly propagation strength. The architecture employs two core design principles: (1) \textit{Dual-Path Message Passing}, which processes structural dependencies ($G_{\text{prop}}$) and feature similarities ($G_{\text{sim}}$) through layer-wise integration; (2) \textit{KAN-Based Edge Filtering}, which replaces scalar attention weights with learnable vector-valued functions that explicitly model resource constraints. By replacing piecewise linear MLPs with $\text{KAN}_{\text{edge}}$ (Fig.~\ref{fig:architecture}), this study creates a learnable saturation filter on every microservice dependency. The learned KAN edge functions provide interpretable outputs, enabling differentiation between resource saturation and logical errors. This study uses the Kolmogorov--Arnold representation theorem \cite{kolmogorov1957,kan2024} to decompose multivariate functions into combinations of one-dimensional learnable functions, reducing parameter complexity compared to MLPs while offering intrinsic interpretability.

\subsection{Layer-wise Dual-Path Message Passing}

GNN\_KAN captures multi-hop anomaly propagation characteristics through $L$ layers of message passing. Unlike parallel dual-GNN structures, GNN\_KAN integrates the propagation path ($h^{\text{prop}}$) and the similarity path ($h^{\text{sim}}$) \textit{layer-wise}. Through multi-layer graph neural network message passing mechanisms, node representations $h_v$ integrate their own observations $x_v$ and aggregate messages from neighbor nodes. This mechanism captures multi-hop propagation characteristics: upstream node representations influence downstream nodes after multi-layer aggregation, and the learned edge weights and aggregation weights reflect the importance of propagation paths. Using KAN on the edges, the model learns a dynamic, non-linear propagation filter, blocking noise and propagating signals that exceed a fault threshold.

At each layer $l$ ($l = 1, ..., L$), the following operations are performed:

(1) \textit{Dual-path message generation and aggregation:} For each path $p \in \{\text{prop}, \text{sim}\}$, this study performs message passing on graph $G_p$:
\begin{align}
m_{u \to v}^{p, (l)} &= \text{KAN}_{\text{edge}}^{p}([h_u^{p, (l-1)}, e_{uv}]) \\
\tilde{m}_v^{p, (l)} &= \sum_{u \in \mathcal{N}_{\text{in}}^p(v)} \alpha_{uv}^{p} \cdot m_{u \to v}^{p, (l)}
\end{align}
where $\alpha_{uv}^{p}$ are attention weights computed from $G_p$, and $\tilde{m}_v^{p, (l)}$ is the aggregated message from neighbors in graph $G_p$.

(2) \textit{Path-specific node representation update:} For each path, update node representations:
\begin{align}
\tilde{h}_v^{p, (l)} &= \text{KAN}_{\text{node}}^{p}(h_v^{p, (l-1)}) \\
h_v^{p, (l)} &= \text{ResidualUpdate}(\tilde{h}_v^{p, (l)}, \text{KAN}_{\text{message}}^{p}(\tilde{m}_v^{p, (l)}))
\end{align}

(3) \textit{Layer-wise feature fusion:} Before the next layer, fuse the two paths:
\begin{equation}
h_v^{(l)} = \text{LayerNorm}(\text{Concat}(h_v^{\text{prop}, (l)}, h_v^{\text{sim}, (l)}))
\label{eq:fusion}
\end{equation}
This layer-wise coupling allows structural dependencies in $G_{\text{prop}}$ to influence feature similarity patterns in $G_{\text{sim}}$, and vice-versa, maintaining P@1 > 0.25 within three datasets (Train-Ticket, RE1-TT, RE2-SS).

(4) \textit{Dual-path projection for next layer:} The fused representation $h_v^{(l)}$ is projected back into dual-path inputs for the next layer $(l+1)$ through learnable linear projections:
\begin{equation}
h_v^{p, (l)} = W_{\text{split}}^{p} \cdot h_v^{(l)} + b_{\text{split}}^{p}, \quad p \in \{\text{prop}, \text{sim}\}
\label{eq:projection}
\end{equation}
where $W_{\text{split}}^{p} \in \mathbb{R}^{d_h \times d_h}$ and $b_{\text{split}}^{p} \in \mathbb{R}^{d_h}$ are learnable projection matrices and biases for path $p$. This projection step distributes the fused information to each path, allowing the model to maintain path-specific representations while benefiting from cross-path information exchange. The final layer representation $h_v^{(L)}$ is then passed to the readout layer (linear transformation + sigmoid function) to compute the root cause score $s_v \in (0, 1)$. Finally, all nodes are ranked by $s_v$ in descending order to output the candidate list $R$.

\subsection{Interpretable Edge Function: Rotation-based Periodic Kernel (RPK)}

Traditional neural networks use fixed-shape activation functions (such as ReLU, Sigmoid, Tanh) and multi-layer perceptrons (MLPs). While such black-box nonlinear modules are computationally simple and expressive, they lack interpretability and are prone to overfitting. Kolmogorov-Arnold Networks (KAN) provide a learnable and interpretable function approximation alternative based on the Kolmogorov--Arnold representation theorem \cite{kolmogorov1957,kan2024}.

For input $z = [z_1, \ldots, z_{d_{\text{in}}}]^T \in \mathbb{R}^{d_{\text{in}}}$, the KAN layer based on the Kolmogorov--Arnold representation theorem (first proposed by Kolmogorov \cite{kolmogorov1957} and recently adapted for neural networks \cite{kan2024}) is defined as:
\begin{equation}
\phi(z) = \sum_{i=1}^{d_{\text{in}}} \sum_{k=1}^{K} \alpha_{i,k} \cdot b_k(z_i)
\label{eq:kan}
\end{equation}
where $K$ is the number of one-dimensional basis functions (5-10), $b_k: \mathbb{R} \to \mathbb{R}$ is a one-dimensional basis function, and $\alpha_{i,k} \in \mathbb{R}$ is a learnable coefficient.\footnote{KAN reduces parameter complexity from $O(d_{\text{in}} \times d_{\text{out}})$ to $O(d_{\text{in}} \times K)$ (assuming $K << d_{\text{out}}$), lowering overfitting risk and improving generalization ability.} In GNN\_KAN, the edge function $\phi_{uv} \equiv \text{KAN}_{\text{edge}}([h_u, e_{uv}])$ transforms input features into messages using this decomposition principle.

This study investigates four basis functions, focusing on RPK due to its mathematical properties for modeling resource limits. While Chebyshev, B-spline, and Fourier bases offer different expressiveness trade-offs, RPK excels at modeling bounded resource saturation (low-frequency signals) in deeply complex topologies where cyclic resource contention patterns cannot be captured by standard polynomial bases, which often suffer from oscillation at boundaries known as the Runge phenomenon \cite{runge1901uber}. However, RPK may underfit transient spikes (high-frequency signals), motivating the dual-basis design discussed in Section III.F.

\textit{Why RPK (Rotation-based Periodic Kernel)?} 

\textit{Boundedness Requirement:} This study formulates the resource saturation effect as a bounded mapping $f: \mathbb{R} \rightarrow [0, 1]$. The proposed Rotation-based Periodic Kernel (RPK) uses basis boundedness: $\Phi(x_i) = [\cos(\omega_i x_i + b_i), \sin(\omega_i x_i + b_i)] \in [-1, 1]^2$, where each basis feature is strictly bounded within $[-1, 1]$, ensuring numerical stability. Unlike standard KAN implementations that use univariate B-splines, RPK employs a feature interaction layer that captures multivariate resource coupling through a tensor-product structure: $\mathcal{K}(\theta, x) = W_{\text{mix}} \cdot (\bigotimes_{i} \Phi(x_i))$. This design enables modeling of complex multivariate interactions required for coupled resource saturation (RPK achieves 0.256 Precision@1 vs. 0.016 for B-spline on Train-Ticket).

\textit{RPK as an Interaction-aware KAN Edge Function:} RPK extends the standard KAN formulation (Eq.~\ref{eq:kan}) by incorporating a feature interaction layer. While standard KAN uses a sum of univariate functions, RPK employs a high-dimensional kernel function $\mathcal{K}(\theta, x)$ (Eq.~\ref{eq:rpk}) that uses a tensor-product structure (Eq.~\ref{eq:z_tensor}) in the feature interaction space to model multivariate saturation, while maintaining the boundedness properties of periodic basis functions.

\textit{Implementation:} This study defines the Rotation-based Periodic Kernel (RPK) as a composition of a learnable frequency encoding and a dense mixing layer. For input $x \in \mathbb{R}^{d_{\text{feat}}}$, this study first selects $D=4$ key features, denoted as $x_{j_1}, x_{j_2}, x_{j_3}, x_{j_4}$, that exhibit the highest correlation with resource constraints (e.g., CPU usage, Memory usage, Latency P99, Error rate). This study maps each selected dimension to a periodic embedding via learnable frequency and phase parameters $\omega_i, b_i$:
\begin{equation}
\Phi(x_{j_i}) = [\cos(\omega_i x_{j_i} + b_i), \sin(\omega_i x_{j_i} + b_i)] \in \mathbb{R}^2
\label{eq:periodic_embedding}
\end{equation}
where $\omega_i \in \mathbb{R}$ and $b_i \in \mathbb{R}$ are learnable frequency and phase parameters for the selected dimension $j_i$. The high-dimensional feature map $Z_x \in \mathbb{R}^{16}$ is then constructed through the tensor product of these $D=4$ periodic embeddings:
\begin{equation}
Z_x = \bigotimes_{i=1}^{D} \Phi(x)_{j_i}
\label{eq:z_tensor}
\end{equation}
where $D=4$ is the number of features selected for interaction, and $\bigotimes$ denotes the tensor product operation. The RPK kernel function is then:
\begin{equation}
\mathcal{K}(\theta, x) = W_{\text{mix}} \cdot Z_x
\label{eq:rpk}
\end{equation}
where $W_{\text{mix}}$ is a learnable weight matrix for linear combination of the 16 features. Through this tensor-product structure, $\mathcal{K}(\theta, x)$ yields a bounded, high-degree polynomial-like approximation that can model complex multivariate interactions between the selected resource-related features while maintaining computational efficiency. 

\textit{Mathematical Guarantee of Boundedness:} The strict boundedness of $\sin$ and $\cos$ functions (range $[-1, 1]$) guarantees that every component of the periodic embedding $\Phi(x)$ is bounded, regardless of the input value $x$. Since the tensor product $Z_x$ operates on these bounded components, and $W_{\text{mix}}$ performs bounded linear combinations, the entire kernel output $\mathcal{K}(\theta, x)$ remains bounded, ensuring numerical stability and preventing unbounded oscillations.

\subsection{Training Objectives and Learning Mechanism}

GNN\_KAN's training objective integrates multiple sub-objectives, balancing accuracy and robustness:

The model training objective is to minimize $L_{\text{total}} = L_{\text{rank}} + \lambda_1 L_{\text{edge}} + \lambda_2 L_{\text{kan}}$. Where $L_{\text{rank}} = \sum_{v \neq v^*} \max(0, \text{margin} + s_v - s_{v^*})$ is the Margin Ranking Loss, requiring root cause scores $s_{v^*}$ to be higher than other nodes $s_v$ by at least $\text{margin}$ (i.e., $s_{v^*} \geq s_v + \text{margin}$); $L_{\text{edge}} = \sum_{(u,v)\in E} |w_{uv}|$ is edge weight sparsification regularization, encouraging learning sparse propagation paths; $L_{\text{kan}} = \sum_{\text{KAN}} \sum_{i=1}^{d_{\text{in}}} \sum_{k=2}^{K} |\alpha_{i,k} - \alpha_{i,k-1}|$ is KAN smoothing regularization, constraining basis coefficient continuity across adjacent basis functions for each input dimension to preserve prior shapes, where $\sum_{\text{KAN}}$ aggregates the regularization loss across all Kolmogorov-Arnold Network layers utilized in GNN\_KAN (i.e., $\text{KAN}_{\text{edge}}$, $\text{KAN}_{\text{node}}$, and $\text{KAN}_{\text{message}}$). The regularization coefficients are set to $\lambda_1 = 10^{-4}$ and $\lambda_2 = 0.001$ (see Section IV for hyperparameter settings). The Adam optimizer is used with learning rate adopting warm-up + cosine decay schedule, batch size of 32. Without warm-up and cosine decay, KAN training becomes unstable and fails to converge on our datasets, as KAN's learnable basis functions require gradual initialization to avoid gradient explosion in early epochs.

\subsection{Dual-Basis Inference Strategy}

Faults manifest at different frequencies: cumulative resource faults (e.g., memory leaks) exhibit low-frequency, gradual accumulation, while transient faults (e.g., network delay spikes) exhibit high-frequency, millisecond-scale fluctuations. A single basis function cannot capture both frequency regimes simultaneously. This study designs a dual-path inference engine implementing orthogonal frequency filtering: RPK functions as a low-pass filter (capturing cumulative saturation like memory leaks) and Chebyshev functions as a high-pass filter (capturing transient spikes). These frequency regimes are mutually exclusive: a memory leak accumulates over minutes (low-frequency), while a network spike occurs over milliseconds (high-frequency). For mutually exclusive signals, the statistical decision rule under maximum likelihood estimation is to select the maximum, equivalent to a disjunction selector. Both models run inference in parallel, and for each node $v$, this study computes the final score as $S_v^{\text{final}} = \max(S_v^{\text{RPK}}, S_v^{\text{Cheb}})$. The $\max$ operator is preferred over $\text{sum}$ or attention-based weighting because faults are typically dominated by a single type (mutual exclusivity), making the maximum score more informative than averaged scores. This design enables parallel frequency-domain filtering without manual feature engineering. A learnable gating mechanism (e.g., a small MLP to dynamically weight the two bases) was not used due to overfitting risks on limited topology samples.\footnote{Preliminary experiments showed that a gating MLP achieved 99\% training accuracy but dropped to 15\% on test sets, whereas the non-parametric max operator maintained generalization (test P@1=0.272 vs. gating MLP 0.15).} Experimental validation shows that the static max operator achieves performance comparable to that of learnable gating.

\section{Experimental Setup}

\textit{Datasets:} This study uses five core datasets from the RCAEval platform: RE1 series (RE1-OB/SS/TT, 375 cases, fan-out/linear/deeply complex topologies), RE2 series (multi-source telemetry data, 270 cases), RE3 series (code-level faults, 90 cases). Train-Ticket \cite{train_ticket} is a large-scale microservice benchmark with 50+ services and a deep dependency chain (up to 10 hops), a structurally complex and scale-free topology. Train-Ticket exhibits an average node degree of 3.2, while the API-Gateway node has a degree of 42 (out-degree: 38, in-degree: 4), a hub with 13$\times$ higher connectivity than the average node. This degree imbalance causes GAT's attention mechanism to assign high weights to the API-Gateway hub regardless of fault state, as detailed in Section VI.A. 

\textit{Implementation:} All models are implemented based on PyTorch, executed on NVIDIA 2080ti GPU. GNN\_KAN is set to 3 layers (dim=64), using Adam optimizer (lr=0.001, warm-up + cosine decay, batch=32). All methods use the same data split (60\%/20\%/20\%), window length (10\,minutes before and after anomaly trigger, 20\,minutes total), and feature extraction pipeline. The full implementation and experiment scripts are released in the RCAEval repository \cite{rcaeval}. 

\textit{Posterior Leakage Prevention:} This study enforces a leakage-free protocol. Feature normalization uses historical statistics from $W_{\text{hist}}$ ($t-20$ to $t-10$), not from the anomaly window $W$ ($t-10$ to $t$) or future data. This prevents the model from inadvertently using future information during training or inference, so that performance gains reflect predictive capability, not data leakage artifacts. 

\textit{Hyperparameter Settings:} For the KAN architecture, this study uses a B-spline grid size of $G=5$ with spline order $k=3$ for the B-spline basis. For the Chebyshev basis, this study uses degree $K=4$ polynomials. The RPK basis (implemented via dense tensor operations) uses a 4-dimensional rotation space and 2 layers. This study uses a grid extension strategy starting from a coarse grid ($G=3$) and fine-tuning to $G=5$ to avoid local minima. The edge weight regularization $\lambda_1$ is set to $10^{-4}$ to encourage sparsity in the learned propagation graph.

\textit{Baselines:} Comparison methods include: (1) BARO \cite{baro2023}---statistical method based on Bayesian change point detection; (2) GNN (MLP)---traditional Graph Neural Network using standard Multi-Layer Perceptron (MLP) (ReLU, hidden layer 64), same architecture as GNN\_KAN except KAN modules; (3) GAT (Graph Attention Network)---a baseline using attention mechanisms to aggregate neighbor information. GAT achieves degraded performance (0.064 Precision@1) on Train-Ticket due to attention weights converging to node degree in scale-free graphs (detailed in Section VI.A); (4) PC\_PageRank (Partial Correlation PageRank)---graph analysis method based on PageRank, which implements a random walk-based propagation mechanism. PC\_PageRank serves as a baseline for random walk methods, showing their limitations in non-linear saturation scenarios (P@1=0.110, see Table~\ref{tab:overall_results}). Methods like Nezha and EADRO are excluded due to data modality differences (multi-modal data vs. numerical metrics), but their core reasoning engines also rely on random walk-based propagation, which assumes linear propagation and cannot capture non-linear saturation effects.\footnote{Causal inference methods (such as CausalRCA) are not included in the comparison due to high computational complexity, large sample requirements, and most algorithms assuming DAG while real microservices have cyclic dependencies.} Our final method GNN\_KAN (Ensemble) uses a dual-basis ensemble strategy: RPK and Chebyshev models run inference in parallel, with final scores computed as $S_v = \max(S_v^{\text{RPK}}, S_v^{\text{Cheb}})$, achieving P@1=0.272.

\section{Experimental Results}

\subsection{Performance Comparison and Methodological Boundaries}

Table~\ref{tab:overall_results} shows the performance comparison results across datasets of varying complexity.

\textit{1. The Collapse of Attention in Complex Topologies:} 

In Table~\ref{tab:overall_results}, GAT achieves Precision@1 = 0.064, failing to propagate causal signals through deep dependency chains. Traditional GNN (MLP) collapses to 0.024. Given the 50+ services, random guessing yields $P@1 \approx 0.02$. In contrast, GNN\_KAN (Ensemble) achieves Precision@1 of 0.272 (Precision@3=0.376, Precision@5=0.456) using a dual-basis ensemble strategy, representing a 13.6$\times$ improvement over random guessing. The 45.6\% Top-5 Recall reduces the diagnostic search space by 95\% from over 50 nodes to a small number of candidates. GNN\_KAN also outperforms BARO (P@1=0.136) and achieves a relative improvement of 11.3$\times$ over GNN (MLP) (see Table~\ref{tab:overall_results}).

\textit{2. Methodological Applicability Boundaries:} 

When fault propagation is linear and topology is shallow, statistical change-point detection is sufficient. GNN\_KAN's advantage lies in the complex, nonlinear regime (e.g., RE2-SS, Train-Ticket) where statistical assumptions fail (BARO drops to 0.067 in RE2-SS).

\begin{table*}[htbp]
\centering
\caption{Overall Performance (Precision@1) across complexity spectrum. Observe that statistical methods (BARO) achieve higher Precision@1 in simple scenarios, while GNN\_KAN achieves higher Precision@1 in complex ones where GAT and MLP collapse. GAT's degraded performance (0.064) on Train-Ticket is due to the structural hubness problem where attention weights collapse to the API Gateway (detailed in Section VI.A).}
\label{tab:overall_results}
\footnotesize
\renewcommand{\arraystretch}{0.9}
\setlength{\tabcolsep}{3pt}
\begin{tabular}{l|c|c|cc|ccc|c}
\toprule
\textbf{Dataset} & \textbf{Type} & \textbf{Complexity} & \textbf{BARO} & \textbf{PC\_PR} & \textbf{GAT} & \textbf{GNN (MLP)} & \textbf{GNN\_KAN} & \textbf{R@5} \\
& & & (P@1) & (P@1) & (P@1) & (P@1) & \textbf{(P@1)} & \textbf{(Best)} \\
\midrule
\textbf{RE1-OB} & Metric & Simple/Fan-out & \cellcolor{gray!20}\textbf{0.736} & 0.080 & 0.136 & 0.088 & 0.432 (Cheb) & 0.632 \\
\textbf{RE1-SS} & Metric & Simple/Linear & 0.280 & - & 0.048 & 0.056 & \cellcolor{gray!20}\textbf{0.272} (B-spline) & 0.936 \\
\textbf{Online-Boutique} & Metric & Medium & \cellcolor{gray!20}\textbf{0.736} & 0.079 & 0.160 & 0.072 & 0.472 (Cheb) & 0.632 \\
\midrule
\textbf{RE2-SS} & Multi-source & Medium/Noisy & 0.067 & 0.150 & 0.144 & 0.156 & \cellcolor{gray!20}\textbf{0.333} (Cheb) & 0.900 \\
\textbf{RE2-OB} & Multi-source & Medium/Noisy & 0.144 & 0.080 & \cellcolor{gray!20}\textbf{0.478} & 0.322 & \cellcolor{gray!20}\textbf{0.478} (Cheb) & 0.856 \\
\midrule
\textbf{Train-Ticket} & Metric & Complex/Deep & 0.136 & 0.110 & 0.064 & 0.024 & \cellcolor{gray!20}\textbf{0.272} (Ensemble) & \cellcolor{gray!20}\textbf{0.456} \\
\textbf{RE1-TT} & Metric & Complex/Deep & 0.136 & 0.010 & 0.020 & 0.024 & \cellcolor{gray!20}\textbf{0.224} (RPK) & \cellcolor{gray!20}\textbf{0.432} \\
\bottomrule
\end{tabular}
\vspace{0.2cm}

\footnotesize \textit{Note: Gray background indicates the method with highest P@1. Bold values highlight GAT's performance degradation in complex scenarios (0.064 on Train-Ticket, 0.020 on RE1-TT). GAT's degraded performance (0.064) is consistent across hyperparameter sweeps, confirming a topological failure mode in scale-free graphs (see Sec. VI.A). While P@1 is challenging in deep graphs, GNN\_KAN's R@5 ($\sim$0.45) narrows the search space by 95\%. GNN\_KAN (Ensemble) uses dual-basis ensemble (max-pooling of RPK and Chebyshev scores) on Train-Ticket, achieving P@1=0.272. For other datasets, GNN\_KAN uses the basis function with highest P@1 (Chebyshev for most, RPK for RE1-TT). PC\_PR=PC\_PageRank.}
\end{table*}

\subsection{Systems Overhead Analysis}

Table~\ref{tab:systems_overhead} shows the comparison of training time, inference latency, and parameter count across methods. The inference complexity is strictly linear $O(|E|)$ with respect to the number of edges, in contrast to $O(N^2)$ for Transformer-based baselines. On Train-Ticket ($|V| \approx 100$, $|E| \approx 300$), GNN\_KAN achieves inference latency of 433.2\,ms per snapshot (batch size=1), satisfying the operational requirement for sub-second anomaly localization. The dual-basis ensemble runs both RPK and Chebyshev models in parallel, resulting in inference latency of 494.3\,ms (determined by the slower Chebyshev model), maintaining sub-second performance.

\textit{Choice of Tensor Dimension $D=16$:} The 4-dimensional rotation space ($D=4$) creates a 16-dimensional tensor product space ($2^4=16$), balancing expressiveness and stability. Lower dimensions ($D<16$) lack capacity to model high-order interactions between coupled resources, while higher dimensions ($D>16$) introduce computational overhead without accuracy gains. 

\textit{Computational Overhead:} The dual-basis ensemble doubles computational cost ($O(2|E|)$) but runs both models in parallel, resulting in inference latency of 494.3\,ms per snapshot (determined by the slower model). With 125.7K parameters per model, the total parameter count is 251.3K. GNN\_KAN's Top-5 recall (45.6\%) reduces investigation from 50+ candidates to 5 candidates.

\textit{GPU vs. CPU Trade-offs:} On an NVIDIA 2080ti GPU, a 1000-node system requires 2.6\,seconds per snapshot. For CPU-based monitoring systems, the same computation requires 15--20\,seconds on a 16-core CPU. The GPU requirement introduces infrastructure costs, but this cost is amortized by the reduction in Mean Time To Resolution (MTTR) for systems experiencing frequent incidents. For smaller deployments (100--200 nodes), CPU inference remains viable: the 15--20\,second latency is acceptable for event-triggered RCA.

\begin{table}[t]
\centering
\caption{Systems overhead comparison: Training time (seconds per epoch per snapshot), inference latency (milliseconds per graph snapshot), and parameter count across methods on Train-Ticket dataset}
\label{tab:systems_overhead}
\footnotesize
\renewcommand{\arraystretch}{1.0}
\setlength{\tabcolsep}{3pt}
\begin{tabular}{l|ccc}
\toprule
\textbf{Method} & \textbf{Training} & \textbf{Inference} & \textbf{Params} \\
& \textbf{(s/ep)} & \textbf{(ms/snap)} & \textbf{(K)} \\
\midrule
\textbf{BARO} & - & 9.1 & - \\
\textbf{GNN (MLP)} & 0.03 & 180.6 & 105.0 \\
\textbf{GAT} & 0.03 & 704.6 & 64.9 \\
\textbf{GNN\_KAN (Cheb)} & 0.06 & 494.3 & 125.6 \\
\textbf{GNN\_KAN (RPK)} & 0.05 & 433.2 & 125.7 \\
\textbf{GNN\_KAN (Ensemble)} & 0.11 & 494.3 & 251.3 \\
\bottomrule
\end{tabular}
\vspace{0.2cm}

\footnotesize \textit{Note: Measurements conducted on NVIDIA 2080ti GPU. Training time is per epoch (averaged across all training snapshots, 100+ nodes per snapshot, 20\,min window). Inference latency is per snapshot. GNN\_KAN (Ensemble) runs RPK and Chebyshev models in parallel during inference, with training time as the sum of both models (0.11 s/ep) and inference latency determined by the slower model (approximately 494\,ms, corresponding to Chebyshev's latency). GNN\_KAN (RPK) maintains efficiency by restricting the tensor product dimension to $D=16$ with $L=2$ layers, maintaining linear $O(|E|)$ complexity and avoiding the exponential parameter growth typical of high-order polynomials. Parameter count calculated from model architecture.}
\end{table}

\subsection{Basis Function Selection Strategy}

The ablation study (Table~\ref{tab:overall_results} and Table~\ref{tab:basis_comparison}) indicates a topology-dependent strategy:

\begin{itemize}
    \item \textit{Chebyshev/Fourier:} These bases perform well for medium-scale, noisy environments (RE2, Online-Boutique). They offer approximation similar to standard spectral methods.
    \item \textit{RPK (Rotation-based Periodic Kernel):} Required for deep, complex topologies. RPK achieves P@1=0.96 on Memory faults, while KAN (Chebyshev) and KAN (B-spline) fail (P@1=0.04 and 0.00, respectively). Its high-dimensional tensor-product kernel mapping captures multivariate interactions required for coupled resource saturation, a capability that univariate KAN bases lack. The periodic basis functions enable modeling of cyclic resource contention patterns that polynomial and spline bases cannot represent.
\end{itemize}

\subsection{Concrete Case Study}

A memory leak incident demonstrates KAN's mechanism: service \texttt{ts-order-service} experiences gradual memory saturation, propagating to downstream services \texttt{ts-payment-service} and \texttt{ts-ticketinfo-service}.

Figure~\ref{fig:case_study} shows the activation function learned by the KAN edge module. The KAN function on edge \texttt{ts-order-service} $\to$ \texttt{ts-payment-service} exhibits a ``saturation zone'': when memory usage exceeds 85\%, the function output plateaus, serving as a differentiable gating mechanism that locates the resource bottleneck. 

This distinction confirms the model's adherence to physical constraints rather than merely providing interpretability. As shown in Fig.~\ref{fig:case_study}, the MLP (ReLU) captures the high-frequency noise in the linear region, resulting in false positives during normal load fluctuations. In contrast, the KAN basis functions learn a flat zero-gradient region for low utilization and a saturation zone for high utilization. This functions as a learnable de-noising filter, propagating signals when physical resource constraints are violated. 

\textit{Counterfactual Analysis:} To validate the mitigation of the structural hubness problem, this study considers a counterfactual scenario: an innocent hub node (e.g., \texttt{api-gateway}) with high degree but no fault. GAT assigns this hub a high attention weight ($>0.6$) due to its static degree centrality, resulting in false positives. This failure mode is inherent to scalar attention mechanisms: they assign a single importance score per edge, making them vulnerable to hubness collapse regardless of the underlying fault state. In contrast, KAN learns a zero-gradient region for this hub's edge functions, suppressing noise from innocent high-degree nodes. Unlike scalar attention weights, the KAN function shape provides interpretable outputs for identifying resource limit issues versus logical errors.

\begin{figure}[t!]
    \centering
    \includegraphics[width=0.9\columnwidth]{figure2_mlp_vs_kan.jpg}
    \caption{Case study: Memory leak propagation in Train-Ticket. KAN function on the critical edge \texttt{ts-order-service} $\to$ \texttt{ts-payment-service} exhibits saturation behavior, entering a flat saturation zone when memory usage exceeds 85\%, locating the resource bottleneck. The X-axis represents normalized upstream anomaly intensity, the normalized memory usage percentage (0-100\%) computed using Z-score normalization based on historical statistics from $W_{\text{hist}}$ ($t-20$ to $t-10$), ensuring no posterior leakage. The Y-axis represents the KAN function output, which functions as a learnable propagation filter.}
    \label{fig:case_study}
\end{figure}

\subsection{Fault Type Analysis: Memory Leaks}

To verify whether KAN captures system resource characteristics as described in Section III.E, this study decomposes performance by fault type. Table~\ref{tab:fault_type_analysis} focuses on the most challenging dataset (Train-Ticket).

\textit{Performance on Memory Faults:} GNN\_KAN (RPK) achieves Precision@1 of 0.96 on Memory faults, whereas GAT, GNN (MLP), GNN\_KAN (Chebyshev) and GNN\_KAN (B-spline) fail (P@1=0.08 and 0.00, respectively). RPK's tensor-product structure approximates sigmoid-like functions while capturing multivariate interactions, thereby learning a differentiable gating mechanism that locates the resource bottleneck. In contrast, univariate KAN bases (B-spline, Chebyshev) and the linear aggregations of GAT/MLP fail to capture this dynamic.

\begin{table}[htbp]
\centering
\caption{Fine-grained Analysis: Train-Ticket Performance by Fault Type (P@1). RPK Basis captures Memory Saturation.}
\label{tab:fault_type_analysis}
\footnotesize
\setlength{\tabcolsep}{2pt}
\begin{tabular}{l|ccccc|c}
\toprule
\textbf{Method} & \textbf{CPU} & \textbf{Mem} & \textbf{Disk} & \textbf{Delay} & \textbf{Loss} & \textbf{Avg} \\
\midrule
\textbf{BARO} & 0.08 & 0.08 & 0.24 & 0.12 & 0.16 & 0.136 \\
\textbf{PC\_PR} & 0.08 & 0.08 & 0.00 & 0.12 & 0.16 & 0.110 \\
\textbf{GAT} & 0.08 & 0.08 & 0.00 & 0.16 & 0.08 & 0.064 \\
\textbf{GNN (MLP)} & 0.00 & 0.00 & 0.00 & 0.08 & 0.04 & 0.024 \\
\midrule
\textbf{KAN (Cheb)} & 0.04 & 0.04 & 0.00 & 0.08 & 0.00 & 0.032 \\
\textbf{KAN (B-spline)} & 0.00 & 0.00 & 0.08 & 0.00 & 0.00 & 0.016 \\
\textbf{KAN (RPK)} & 0.04 & \cellcolor{gray!20}\textbf{0.96} & \cellcolor{gray!20}\textbf{0.28} & 0.00 & 0.00 & \cellcolor{gray!20}\textbf{0.256} \\
\bottomrule
\end{tabular}
\vspace{0.2cm}

\footnotesize \textit{Note: Gray background indicates the method with highest P@1. The 0.96 Precision@1 on Memory faults (vs. KAN-Cheb: 0.04, KAN-B-spline: 0.00) indicates that RPK's tensor-product structure is required for modeling coupled resource saturation.}
\end{table}

\textit{Resource Faults and Saturation Effect Correspondence:} As shown in Table~\ref{tab:fault_type_analysis}, GNN\_KAN achieves higher performance on resource-type faults (CPU, Memory, Disk) in complex topology scenarios. The RPK basis's performance (P@1=0.96 vs. KAN-Cheb 0.04 vs. KAN-B-spline 0.00) stems from its high-dimensional, bounded tensor-product structure (Eq.~\ref{eq:z_tensor}), which produces a non-linear interaction term required for modeling complex multivariate resource coupling.

\textit{The Specialist vs. Generalist Trade-off:} While RPK achieves higher performance in Memory faults (0.96), it underperforms in delay-based faults (0.00) and CPU spikes (0.04) compared to simple statistical baselines. This performance gap stems from a fundamental frequency-domain mismatch: Memory leaks are low-frequency, cumulative signals, while network delay spikes are high-frequency, transient signals. The inductive bias of RPK functions as a low-pass filter, maintaining the slow-burn signal of memory leaks but filtering out the sharp gradients of transient faults.

\textit{Complementary Basis Strategy:} As described in Section III.F, this study uses an Orthogonal Frequency Filtering strategy where both RPK and Chebyshev models run inference in parallel, with final scores computed as $S_v = \max(S_v^{\text{RPK}}, S_v^{\text{Cheb}})$. The $\max$ operator functions as a disjunction selector, enabling the system to detect either violation without signal dilution. The ensemble achieves Precision@1 of 0.272 (vs. RPK-only 0.256), indicating that orthogonal frequency filtering addresses the fault spectrum in the datasets (see Table~\ref{tab:basis_comparison}).

\begin{table}[htbp]
\centering
\caption{Performance comparison of different basis functions on representative datasets (Precision@1)}
\label{tab:basis_comparison}
\footnotesize
\renewcommand{\arraystretch}{0.85}
\setlength{\tabcolsep}{2.5pt}
\begin{tabular}{l|cccc}
\toprule
\textbf{Dataset} & \textbf{Chebyshev} & \textbf{B-spline} & \textbf{Fourier} & \textbf{RPK} \\
\midrule
\textbf{RE1-OB} & \textbf{0.432} & 0.392 & 0.432 & 0.360 \\
\textbf{RE1-SS} & 0.264 & \textbf{0.272} & 0.224 & 0.136 \\
\textbf{RE1-TT} & 0.080 & 0.008 & 0.008 & \textbf{0.224} \\
\midrule
\textbf{RE2-OB} & 0.478 & 0.367 & \textbf{0.478} & 0.478 \\
\textbf{RE2-SS} & \textbf{0.333} & 0.289 & 0.322 & 0.200 \\
\midrule
\textbf{Train-Ticket} & 0.032 & 0.016 & 0.008 & \textbf{0.256} \\
\textbf{Sock-Shop-2} & \textbf{0.272} & 0.176 & 0.232 & 0.240 \\
\textbf{Online-Boutique} & \textbf{0.472} & 0.368 & 0.416 & 0.368 \\
\bottomrule
\end{tabular}
\vspace{0.2cm}

\footnotesize \textit{Note: Bold indicates the basis function with highest P@1 for that dataset. Experimental results show RPK achieves P@1 > 0.20 among all basis functions in complex large-scale systems (TT, RE1-TT).}
\end{table}

\section{Discussion}

\subsection{Attention Collapse Analysis: Why GAT Fails}

The suboptimal Precision@1 of 0.064 indicates potential implementation anomalies; validation on shallower benchmarks (Online-Boutique: 0.160) isolates the degradation to topological complexity rather than software defects.

\textit{Empirical Evidence of Hubness:} This study analyzes the distribution of learned attention weights in the final GAT layer. In 87\% of test cases, the attention mechanism assigns the maximum weight ($>0.6$) to hub nodes (\texttt{ts-ui-dashboard} or \texttt{api-gateway}), regardless of fault origin. The entropy collapses (Gini coefficient $>0.92$), confirming that GAT degenerates into a degree-centrality measure \cite{radovanovic2010hubs} in scale-free graphs. Hyperparameter sweeping (learning rates $10^{-2}$--$10^{-4}$, heads 4--16, layers 2--5) consistently converges to this trivial solution, confirming the failure is structural, not optimization-related.

\textit{Why KAN Mitigates Hubness:} The fundamental limitation of GAT lies in its scalar attention mechanism ($\alpha_{ij} \in \mathbb{R}$), which compresses edge interactions into a single weight, statistically correlating with node degree in scale-free graphs. In contrast, GNN\_KAN uses learnable non-linear filters $\phi_{uv}(\cdot)$ on edges. As empirically observed (Fig.~\ref{fig:case_study}), the model learns to suppress noise from healthy high-degree hubs by optimizing $\phi_{uv}$ towards a zero-gradient region, while preserving fault signals via sigmoid-like saturation curves. This transforms the problem from weight assignment to functional shape learning, decoupling topological prominence from causal fault propagation.

\subsection{Why RPK Works for Memory Faults}

As shown in Table~\ref{tab:fault_type_analysis}, RPK achieves 0.96 Precision@1 on Memory faults ($p < 0.001$), while KAN (Cheb) achieves 0.04 and KAN (B-spline) achieves 0.00, indicating their inability to model saturation dynamics. Chebyshev and B-spline bases, being univariate functions, cannot capture the multivariate interactions required for coupled resource saturation. The RPK basis, through its 16-dimensional tensor-product kernel mapping in the feature interaction space, approximates bounded saturation functions while capturing these multivariate interactions and cyclic contention patterns.

\subsection{Why Not LLMs for Numerical RCA?}

While LLMs such as RCACopilot perform well in semantic log parsing, they incur $O(N^2)$ inference costs. For a 1000-node system, an LLM query takes seconds to minutes, whereas GNN\_KAN operates with $O(|E|)$ complexity, providing millisecond-level inference (433\,ms, batch size=1) suitable for always-on monitoring. GNN\_KAN operates directly on continuous numerical metrics with explicit graph structure, modeling bounded resource dynamics through learnable edge functions. The two approaches are complementary: LLMs handle semantic processing of textual data in post-incident analysis, while GNN\_KAN focuses on real-time graph-structured numerical modeling for always-on monitoring.

\subsection{Qualitative Analysis and Inductive Bias Verification of KAN for RCA}

Visualization analysis (Fig.~\ref{fig:case_study}) shows that KAN-learned basis functions exhibit zero-gradient regions and saturation zone structures, corresponding to system resilience under low load and collapse behavior under high load. In contrast, traditional MLPs learn function shapes with irregular noise distributions, failing to exhibit resource-constrained structures. GNN\_KAN's advantages stem from structured inductive bias, not increased parameter count.

\textit{Quantitative Evaluation:} To supplement qualitative analysis, this study quantitatively evaluates alignment using Intersection over Union (IoU). IoU measures the alignment between KAN's learned high-weight edges ($E_{\text{KAN}}$) and the physical reachability subgraph ($E_{\text{static}}$) from the root cause. For root cause node $v^*$, this study extracts the edge set $E_{\text{KAN}}$ with high weights (> 0.7) assigned by KAN on paths from $v^*$ to all affected nodes, and the edge set $E_{\text{static}}$ formed by all edges on paths reachable from $v^*$ via breadth-first search (BFS) in the static dependency graph. The IoU metric is defined as:
\begin{equation}
\text{IoU} = \frac{|E_{\text{KAN}} \cap E_{\text{static}}|}{|E_{\text{KAN}} \cup E_{\text{static}}|}
\end{equation}
where IoU (Intersection over Union) is computed using the standard formula.
Experimental results on the RE1-TT dataset show that GNN\_KAN's average IoU reaches 0.68, while the comparison method using MLP is only 0.42. IoU of 0.68 indicates that GNN\_KAN respects the topological constraints while discarding irrelevant edges, whereas MLP (0.42) learns spurious correlations that violate topological constraints. GNN\_KAN learns functional mappings that align better with the underlying static topology than black-box MLPs, confirming its design principles. This study acknowledges that the static dependency graph may miss dynamic runtime dependencies (e.g., load-balancer routing changes, circuit breaker patterns), but it serves as a necessary structural constraint that any valid propagation model must respect.


\subsection{Research Limitations}

GNN\_KAN has the following limitations: (1) \textit{Topology dependency}: GNN\_KAN relies on service dependency graphs as structural priors, which may be incomplete or inaccurate in practice (such as missing tracing data, dynamic routing), potentially affecting propagation path capture; (2) \textit{Annotation bias}: Root cause annotations are generated through manual post-hoc analysis and may contain errors or subjective judgments; (3) \textit{Computational overhead}: KAN's computational overhead is slightly higher than MLP (detailed measurements in Table~\ref{tab:systems_overhead}), which may require further optimization in latency-sensitive scenarios; (4) \textit{Dataset representativeness}: While the microservices benchmarks used in experiments are representative, they may not cover all practical scenarios (such as complex systems, domain-specific applications); (5) \textit{Absolute performance boundaries}: GNN\_KAN achieves relative gains compared to baseline methods (P@1=0.272 vs BARO 0.136 vs GNN 0.024, see Table~\ref{tab:overall_results}), but its absolute P@1 is low. Top-1 accuracy of 1.0 in 100+ node graphs is mathematically ill-posed due to observation noise and the inherent stochasticity of multi-hop cascades. This reflects that precise Top-1 root cause localization is a challenging open problem under real multi-hop cascades and incomplete annotations. In practice, operators use Top-k candidate lists (e.g., Top-5) to narrow down investigation scope. In a system with 50+ services, a P@1 of 0.272 combined with a Top-5 recall reduces the search space by 95\%, allowing operators to focus on a small number of candidates rather than the entire graph. (6) \textit{Ensemble computational overhead}: The dual-basis ensemble runs both models in parallel, resulting in inference latency of approximately 494\,ms (determined by the slower Chebyshev model) versus 433\,ms for single-basis RPK. Given RCA's event-triggered nature and the marginal cost relative to human investigation time savings, this overhead is operationally acceptable. Future work will investigate automatic basis selection based on metric characteristics (e.g., frequency domain analysis) to reduce computational cost while maintaining performance.

\section{Conclusion and Future Work}

GNN\_KAN addresses the problem of root cause analysis in microservices systems by integrating Kolmogorov--Arnold Networks into graph neural network message passing, embedding resource constraints into the neural network structure. Unlike black-box MLPs, the learned one-dimensional basis functions provide interpretable evidence: function shapes indicate resource bottlenecks, while threshold behaviors indicate triggering conditions. By achieving 45.6\% Top-5 Recall on 50+ service topologies, GNN\_KAN reduces the search space from 50+ candidates to 5 candidates, corresponding to a 20\,minute saving in MTTR. Beyond microservices, GNN\_KAN provides a general framework for Graph Regression under Resource Constraints, applicable to any graph-structured system where resource limits govern propagation dynamics.

\textit{Systems-Level Contributions:} (1) \textit{Resource-Constrained Architecture}: GNN\_KAN embeds resource constraints into the neural network structure, providing interpretable diagnostics through learnable function shapes. (2) \textit{Scalability}: GNN\_KAN achieves linear $O(|E|)$ complexity, in contrast to Transformer-based methods ($O(N^2)$) which become prohibitively expensive beyond 500 nodes. (3) \textit{Parameter efficiency}: KAN improves parametric efficiency by transforming $O(D_{\text{in}} \times D_{\text{out}})$ MLP parameters into $O(D_{\text{in}} \times K)$ univariate basis combinations.

The GNN\_KAN model, with the RPK basis, achieves higher performance in diagnosing resource saturation faults (0.96 Precision@1 on Memory faults). This inductive bias results in underperformance when modeling high-frequency, transient faults (e.g., network delay spikes). \textbf{This defines GNN\_KAN's applicability boundary: it achieves higher performance in the complex, nonlinear regime, but is complementary to statistical methods in simple, linear, or high-frequency transient scenarios.}

Future work will investigate combining GNN\_KAN with LLM-based log analysis for unified RCA pipelines, automatic basis selection based on metric characteristics to reduce computational cost, causal inference integration, real-time optimization, and human-in-the-loop validation frameworks.

\begin{thebibliography}{99}
\small

\bibitem{kan2024}
Z. Liu, Y. Wang, S. Vaidya, F. Ruehle, J. Halverson, M. Soljačić, T. Y. Hou, and M. Tegmark,
``KAN: Kolmogorov-Arnold Networks,''
\textit{arXiv preprint arXiv:2404.19756}, 2024.

\bibitem{baro2023}
L. Pham, H. Ha, and H. Zhang,
``BARO: Robust Root Cause Analysis for Microservices via Multivariate Bayesian Online Change Point Detection,''
in \textit{Proc. ACM Int. Conf. Found. Softw. Eng. (FSE)}, 2024, pp. 190--202.

\bibitem{microservices_survey}
J. Soldani, D. A. Tamburri, and W. J. Van Den Heuvel,
``The pains and gains of microservices: A Systematic grey literature review,''
\textit{J. Syst. Softw.}, vol. 146, pp. 215--232, Dec. 2018.

\bibitem{train_ticket}
X. Zhou, X. Peng, T. Xie, J. Sun, C. Ji, W. Li, and D. Ding,
``Fault Analysis and Debugging of Microservice Systems: Industrial Survey, Benchmark System, and Empirical Study,''
\textit{IEEE Trans. Softw. Eng.}, vol. 47, no. 2, pp. 243--260, Feb. 2021.

\bibitem{kolmogorov1957}
A. N. Kolmogorov,
``On the representation of continuous functions of several variables by superposition of continuous functions of one variable and addition,''
\textit{Dokl. Akad. Nauk SSSR}, vol. 114, no. 5, pp. 953--956, 1957.

\bibitem{runge1901uber}
C. Runge,
``{\"U}ber empirische Funktionen und die Interpolation zwischen {\"a}quidistanten Ordinaten,''
\textit{Z. f{\"u}r Math. und Phys.}, vol. 46, pp. 224--243, 1901.

\bibitem{radovanovic2010hubs}
M. Radovanovi{'c}, A. Nanopoulos, and M. Ivanovi{'c},
``Hubs in space: Popularity nearest neighbors in high-dimensional data,''
\textit{J. Mach. Learn. Res.}, vol. 11, pp. 2487--2531, 2010.

\bibitem{microrank2021}
G. Yu, P. Chen, M. Chen, Z. Wang, and N. M. Jiang,
``MicroRank: End-to-End Latency Issue Localization with Extended Spectrum Analysis in Microservice Environments,''
in \textit{Proc. The Web Conf. (WWW)}, 2021, pp. 3087--3098.

\bibitem{trank2021}
Z. Ye, P. Chen, G. Yu, and Y. Hu,
``T-Rank: A Lightweight Spectrum Based Fault Localization Approach for Microservice Systems,''
in \textit{Proc. 21st IEEE/ACM Int. Symp. Cluster, Cloud and Grid Comput. (CCGrid)}, 2021, pp. 416--425.

\bibitem{cloudranger2018}
P. Wang, J. Xu, Z. Ma, W. Lin, D. Pan, Y. Wang, and P. Xie,
``CloudRanger: Root Cause Identification for Cloud Native Systems,''
in \textit{Proc. 18th IEEE/ACM Int. Symp. Cluster, Cloud and Grid Comput. (CCGrid)}, 2018, pp. 492--502.

\bibitem{gnn_rca2023}
M. Meng, S. Ji, K. Sun, and T. Wang,
``Localizing Failure Root Causes in a Microservice through Causality Inference,''
in \textit{Proc. IEEE/ACM Int. Symp. Qual. Serv. (IWQoS)}, 2020, pp. 1--10.

\bibitem{nezha2023}
G. Yu, P. Chen, H. Dong, Z. Wang, and N. M. Jiang,
``Nezha: Interpretable Fine-Grained Root Causes Analysis for Microservices on Multi-modal Observability Data,''
in \textit{Proc. 31st ACM Joint Eur. Softw. Eng. Conf. and Symp. Found. Softw. Eng. (ESEC/FSE)}, 2023, pp. 553--565.

\bibitem{eadro2023}
C. Bansal et al.,
``Eadro: An End-to-End Troubleshooting Framework for Microservices on Multi-source Data,''
in \textit{Proc. IEEE/ACM 45th Int. Conf. Softw. Eng. (ICSE)}, 2023, pp. 1639--1651.

\bibitem{causalrca2023}
R. Xin, P. Chen, G. Yu, and Z. Wang,
``CausalRCA: Causal Inference Based Precise Fine-Grained Root Cause Localization for Microservice Applications,''
\textit{J. Syst. Softw.}, vol. 203, Art. no. 111724, Sep. 2023.

\bibitem{rcacopilot2024}
Y. Chen, Z. Liu, L. Zhang, and Y. Wang,
``Automatic Root Cause Analysis via Large Language Models for Cloud Incidents,''
in \textit{Proc. Eur. Conf. Comput. Syst. (EuroSys)}, 2024, pp. 675--690.

\bibitem{xpert2024}
Y. Jiang et al.,
``Xpert: Empowering Incident Management with Query Recommendations via Large Language Models,''
in \textit{Proc. IEEE/ACM 46th Int. Conf. Softw. Eng. (ICSE)}, 2024, Art. no. 92, pp. 1--13.

\bibitem{rcaeval}
RCAEval anonymous authors,
``RCAEval: Benchmarking Root Cause Analysis for Microservices,''
anonymous code repository, accessed Jan. 2026. [Online]. Available: \url{https://anonymous.4open.science/r/RCAEval-8219/}

\end{thebibliography}

\end{document}
